\documentclass[10pt,a4paper]{amsart}
\usepackage[margin=19mm]{geometry}
\usepackage{amsmath,amssymb,mathtools}
\usepackage[scr=rsfs,cal=euler]{mathalfa}
\usepackage{xfrac}
\usepackage[unicode,hidelinks,bookmarks=false]{hyperref}
\newcommand{\ie}{i.e.\ }
\newcommand{\cf}{cf.\ }
\newcommand{\resp}{respectively\ }
\newcommand{\Subsection}{\subsection}
\providecommand{\nobreakdash}{\nobreak\hbox{-}}
\setlength{\emergencystretch}{2em}
\multlinegap0pt
\usepackage{tikz}
\usepackage{amssymb}
\usepackage{xcolor}
\usepackage{algorithm}\usepackage{algpseudocode}
\usepackage{float}
\usepackage{dsfont}
\usepackage[shortlabels]{enumitem}
\newtheorem{proposition}{Proposition}
\title{Burning rooted graph products}
\author{John Peca-Medlin}
\date{Source: arXiv:2603.00304v1 (CC BY 4.0)}
\begin{document}
\maketitle

\noindent\emph{Source and licence.} Burning rooted graph products. Version arXiv:2603.00304v1 (CC BY 4.0), distributed by arXiv under the Creative Commons Attribution 4.0 International licence, http://creativecommons.org/licenses/by/4.0/, as recorded in the arXiv licence metadata for this exact version and shown on the abstract page https://arxiv.org/abs/2603.00304v1. Full licence notice: \texttt{licenses/arxiv-2603.00304-CC-BY-4.0.txt}.\par\smallskip

\noindent\textit{Editorial scope.} Counted unit: the article's proposition normalize\_greedy\_comb (source0.tex lines 1644--1649). The article's complete original proof of it is delivered with it (source0.tex lines 1651--1777, one \texttt{proof} environment of 127 lines). No mathematics is written, repaired or re-derived: the statement and the proof are the article's own lines, in the article's own notation, and no retained line is substituted or re-flowed.  No further source-local context is needed: every cross-reference of every retained line resolves inside the two delivered spans.  Nothing was omitted from any retained span, so the package carries no omission record.  No retained line contains a citation, so the wrapper carries no bibliography and no bibliography entry.  No retained line contains a figure environment, an image-inclusion command or drawing code, so no image is delivered and the package declares no asset dependency.  Editorial formatting only, in addition: an AMS \texttt{amsart} preamble that restates the article's own declarations and macros used by the retained lines, namely the environments \texttt{proposition} on separate counters, together with the standard packages those lines need.  The retained environments print in this document as Proposition 1 in order of appearance, whereas the article prints the same items with the numbering of its own full document.  The complete native source of the article as distributed in the arXiv e-print for this exact version is bundled unchanged as \texttt{sources/195479/source0.tex}.
\begin{proposition}[Normalization to \textsc{GreedyComb}]\label{prop: normalize_greedy_comb}
Let $n \ge m$ and suppose $\mathcal B$ is a successful burning sequence
that burns $C_{n,m}$ in $T$ rounds.
Then $\mathcal B$ can be transformed, without loss of coverage,
into the burning sequence generated by $\textsc{GreedyComb}(T;n,m;1)$.
\end{proposition}

\begin{proof}
Label vertices as $(i,j)$,
where $j\in\{1,\dots,n\}$ indexes the tooth
and $i\in\{1,\dots,m\}$ indexes height along the tooth,
with $(1,j)$ the spine vertex
and $(m,j)$ the leaf. Label the burning sequence $\mathcal B=\{(v_t,r_t)\}_{t=1}^T$
with burning radii $r_t=T-t$.

Our goal is to describe an explicit sequence of steps that transforms this initial burning sequence into the one produced by \textsc{GreedyComb}$(T;n,m;1)$. By construction, that sequence is determined by placing the first $T-m+1$ fires along the spine, maximally spaced to burn the left-most teeth from left to right, and then placing the final $m-1$ fires at leaf vertices in the lower right of the comb, proceeding from right to left to eliminate any remaining uncovered segments.

\medskip
\noindent
\emph{Step 0: Reduction.}
If some fire $(v_s,r_s)$ burns no vertex
not already burned by earlier fires,
remove it and increase the radii of all later fires by $1$.
This corresponds to placing all later fires one round earlier (so burning still completes after $T$ rounds). This will be repeated after each following step.

\medskip
\noindent
\emph{Step 1: Move large fires to the spine.}
If $r_t \ge m-1$,
move the fire to the spine vertex $(1,j)$ of its tooth.
Since the leaf $(m,j)$ is distance $m-1$ from the spine,
this preserves vertical reach and weakly increases horizontal coverage.
After each such move, apply Step~0 if needed.
Hence all fires with radius at least $m-1$ lie on spine vertices; these are precisely the first $T-m+1$ fires.

\medskip
\noindent
\emph{Step 2: Boundary tightening.}
A spine fire at tooth $j$ completely burns every tooth
within horizontal distance $r_t-(m-1)$ of $j$.
If such a fire extends beyond tooth $1$ or $n$,
shift it inward so that one extremal burned tooth
is exactly $1$ or $n$.
If this is impossible,
then that fire burns all teeth of the comb.
In this case, all later fires are redundant by Step~0,
and the configuration reduces to a single spine fire
whose placement may be chosen to coincide
with the first placement of \textsc{GreedyComb}.

\medskip
\noindent
\emph{Step 3: Spreading spine fires.}
Suppose two consecutive spine fires both completely burn the same tooth. 
Shift the left fire as far left as possible and, if necessary, shift the right fire as far right as possible so that no tooth is fully burned by both fires. If such a separation is impossible, then their union already burns the entire comb. In this case, move the left fire so that its leftmost fully burned tooth is the first tooth, and then shift the right fire as far right as possible. 

After finitely many such adjustments, the collections of fully burned teeth
corresponding to distinct spine fires are pairwise disjoint (or the entire comb is consumed by spine fires).
We may then reorder the spine fires so that earlier fires lie to the left of later fires.
Each spine fire induces a trapezoidal burning region: a contiguous block of fully burned teeth together with triangular boundary portions on adjacent teeth.
When two such regions abut along their fully burned blocks,
their union is again a trapezoid.
So this reordering preserves the total union of burned vertices and does not change overall coverage.

The spine fires therefore form one or more disjoint blocks
separated by triangular uncovered regions along the leaves of teeth not fully consumed by a spine fire. If the spine fires  already consume the entire graph, then go to Step 6.

\medskip
\noindent
\emph{Step 4: Counting uncovered leaves.}
Let $n'$ be the number of remaining fires (after reduction).
Since each remaining fire can burn at most one leaf,
and the input burning sequence succeeds,
the number of uncovered leaves after the spine fires
is at most $n' \le m-1$.
Moreover, these remaining fires have radii at most
$m-2,m-3,\dots,1,0$.

\medskip
\noindent
\emph{Step 5: Greedy reassignment on uncovered segments.}
Each uncovered region forms a vertical segment
in a tooth adjacent to a boundary of spine coverage.
Assign the remaining $n'$ fires greedily:
place the earliest remaining fire
on the leaf of the largest uncovered segment,
breaking ties by choosing the right-most such segment.
Proceed inductively.

Since $n'$ bounds the number of uncovered leaves,
this assignment burns all remaining segments,
using exactly one fire per segment.
Thus each tooth contains at most one fire,
and every fire is either a spine fire or a leaf fire.

\medskip
\noindent
\emph{Step 6: Left-shift normalization.}
If some leaf fire lies strictly to the left of another,
let $\ell$ be the leftmost such leaf fire.
Shift every fire strictly to the right of $\ell$
one tooth to the left.

Because the spine fires collectively burn
a fixed number of fully burned teeth,
this shift preserves that number.
The rightmost boundary of the spine coverage
has triangular slope $-1$,
so shifting left by one tooth removes the leftmost uncovered leaf
and creates exactly one new uncovered leaf at the right boundary.
Hence the number of leaves unburned by spine fires is preserved.

Reassign the leaf fires greedily as in Step~5.
Coverage by time $T$ is preserved,
since fully burned teeth are unchanged
and only boundary leaves are reallocated. 

Each shift strictly decreases
the number of leaf fires lying left of the rightmost block.
Since there are at most $m-1$ leaf fires,
this process terminates after finitely many steps.

\medskip
At termination,
all leaves not burned by spine fires
lie on the rightmost teeth.
The spine fires occupy consecutive spine vertices,
ordered by decreasing radius from left to right.
The remaining leaf fires occupy the rightmost teeth,
ordered by decreasing radius from right to left.

This configuration is precisely the burning sequence produced
\textsc{GreedyComb}$(T;n,m;1)$.
\end{proof}

\end{document}
